% ============ Criterio: entendimiento ============
\section{Entendimiento del objetivo del repositorio}

El objetivo del repositorio es claro y bien documentado: consolidar, validar y analizar el padrón de investigadores reconocidos por MinCiencias y producir un observatorio crítico con evidencia sobre la fiabilidad, comparabilidad y equidad del sistema de reconocimiento. La documentación (README, informe\_final.md, catálogo YAML) permite entender sin ambigüedad qué se buscaba. Se observa una evolución del alcance: el informe inicial trabajaba con 3 convocatorias (2017/2019/2021) y la versión final del README cubre 6 (2013-2021), incorporando producción, diversidad y redes; esta ampliación es coherente con un observatorio en maduración, aunque obliga a reinterpretar las métricas de las primeras fases. Punto de mejora: formalizar un marco de indicadores (definición operativa, fórmula, fuente y límites de cada métrica) para que la entrega sea auditable.


\subsection{Preguntas de investigación que el repositorio responde}

El repositorio declara 7 preguntas (README): fiabilidad de los datos,
retención entre convocatorias, concentración territorial, brecha de género,
redes de co-filiación, diversidad poblacional y productividad cruzada con el
dataset de producción. Cada pregunta mapea a un módulo de análisis
(\texttt{src/analisis/*}), lo que demuestra una descomposición
objetivo $\rightarrow$ preguntas $\rightarrow$ módulos correcta y trazable.
